In contrast to sequential approaches, interleaved methods cast motion planning as a search problem, alternating between state space sampling and control propagation to find a valid trajectory.
These methods often employ steering functions, or inverse models, to delineate a set of controls connecting two arbitrary points in the state space. 
For instance, DIMT-RRT utilizes a quadratic steering function to account for joint acceleration limits and non-zero initial and goal velocities \cite{kunz2014probabilistically}.
AVP computes the range of reachable velocities at the end of a path, given reachable velocities at the start of it \cite{pham2013kinodynamic,johnson2012optimal}.
LQR-RRT* uses a quadratic cost function and linearized dynamics in conjunction with a linear-quadratic regulator (LQR) to connect consecutive states in an RRT* tree \cite{perez2012lqr,caldwell2018fast}.
Notably, the optimization techniques discussed in sequential methods can also be used as steering functions in interleaved approaches.
In all, these techniques enable precise goal state attainment, allow to perform backward searches, and facilitate exact tree connections in bidirectional searches, thus expediting the planning process \cite{lavalle2006planning}.
However, they come with a set of challenges. Steering functions may not exist or could be difficult to design for a given dynamical system. Furthermore, exact state connections may not be reliable in real-world settings due to uncertainties and the use of optimization as steering may cause planning to get stuck in local minima \cite{littlefield2020efficient}.
